\section{Evaluación General}
\label{sec:evaluacion}

Se integran aquí las métricas de las tres líneas de análisis desarrolladas, evaluadas siempre sobre datos no vistos durante el entrenamiento (conjuntos de prueba o semanas futuras).

\subsection*{Aprendizaje supervisado}

El modelo seleccionado (Regresión Lineal) alcanzó un R\textsuperscript{2} de 0.9992 y un RMSE de 6.94 W en el conjunto de prueba (Sección \ref{sec:supervisado}, Tabla de comparación de modelos).

\subsection*{Aprendizaje no supervisado}

El agrupamiento de los 24 inversores en 3 clústeres y la detección de anomalías (1.99\% de los registros diurnos) proporcionan un mecanismo de diagnóstico complementario, sin necesidad de datos etiquetados de fallas (Sección \ref{sec:no_supervisado}).

\subsection*{Forecasting}

El modelo Holt-Winters superó al baseline ingenuo estacional con un MAPE de 21.79\% frente a 29.64\% (Sección \ref{sec:forecasting}), confirmando que la estacionalidad anual de la generación solar es explotable para la planeación a mediano plazo.

\subsection*{Síntesis}

\begin{table}[H]
\centering
\begin{tabular}{lll}
\toprule
\textbf{Línea de análisis} & \textbf{Métrica principal} & \textbf{Resultado} \\
\midrule
Supervisado (predicción de potencia) & R\textsuperscript{2} / RMSE & 0.9992 / 6.94 W \\
No supervisado (diagnóstico) & \% registros anómalos & 1.99\% \\
Forecasting (energía semanal) & MAPE & 21.79\% \\
\bottomrule
\end{tabular}
\caption{Resumen de evaluación de las tres líneas de análisis}
\end{table}
